        \btabularx
        \hline
        \multicolumn{4}{|l|}{\bf static mode attributes} \\
        \hline
        init       & 3PL  &           & initial value             \\
        domain     & 3PL  & clock     & force the clock domain    \\
        continuous & 3PL  & boolean   & no startup dependence     \\
        readonly   & 3PL  & boolean   & set to read only          \\
        alu        & 3PL  & boolean   & use an ALU if appropriate \\
        dsp        & 3PL  & boolean   & use a DSP block           \\
        loc        & UC:C & string(s) & locations                 \\
        iob        & UC:C & boolean   & locate in I/O blocks      \\
        ids        & UC:C & string(s) & element IDs               \\
        \hline
        \btabularxend
        
        3PL - 3PL uses the attribute itself for control logic or optimisation.
        
        UC:C - 3PL does not use the attribute - it must be used by 3PL source
        code to generate constraint file entries.
        
        {\bf These attributes are case-insensitive.}
        
        {\bf init} sets an initial value for the variable. Any previous initial
        value will be overwritten. If the static variable type is compound the
        initialisation value will usually be a compound value, but it may be
        initialised by a single immediate value in which case all members of the
        compound type will be initialised to that value (the compound type
        members would have to be the same primitive type). The initial value
        is not used until completion of immediate execution when variables
        are processed into target code and hence it can be changed at any time
        during immediate execution.

        {\bf domain} forces the clock domain (both input and output) to
        the specified clock domain. If assigned {\bf null} the current
        clock domain is selected. If the variable already has a clock domain
        as a result of having been assigned or evaluated then that clock domain
        is compared with the intended domain. If the assigned clock domain
        agrees with that already in place then there is no further action,
        otherwise a fatal error occurs.

        {\bf Continuous} true specifies that if possible assignment will occur
        as soon as the FPGA is operational after configuration. This can
        increase operating speed by eliminating clock-enable logic. There are
        two situations where this can occur -
        \blist
        \item   the static variable has a non-conditional assignment - the
                assignment operates continuously following configuration loading
                regardless of any 3PL startup mechanisms.
        \item   a {\bf when} statement has only unconditional assignments to
                variables in its nested code block(s) - the assignments are
                enabled directly from the {\bf when} test expression (or its
                inverse) and are not controlled from execution threads
                initiated by the 3PL startup mechanisms.
        \blend
        see {\bf start()}(\autorefph{sec:ipstart}).
        Note that if directive "continuous" is defined as type log with value
        true, all static variables which do
        not have this attribute set will be treated as though they had the
        continuous attribute true. The directive prevails if the attribute has
        not been set, otherwise the attribute is used.
        
        The "continuous" attribute should be set prior to assignments to the
        variable to have effect. Similarly the "continuous" directive should
        be set as early as possible in immediate execution or else prior
        to immediate execution using presetdirective().

        {\bf readonly} true specifies that from now on this variable cannot
        be assigned to. This attribute cannot be assigned false.
        
        {\bf alu} specifies that assignments of expressions containing
        arithmetic operators '+' and '-' as the final or only operation can be
        combined into a single common adder/subtracter using CLBs. If the {\bf alu}
        attribute is false it can override the global directive {\bf ALU} to
        suppress the use of an ALU for an individual static variable.
        
        The {\bf alu} attribute can save logic and incurs no time penalty for the
        operations drawn in to the adder/subtracter, but all other assignments
        suffer an increase in propagation time since assigned values must now
        traverse the shared adder/subtracter. This tradeoff must be considered
        in each case.
        
        The {\bf alu} attribute can only be applied to a static variable of
        primitive type; it will be ignored following a warning message
        otherwise.

        The {\bf alu} attribute is only used at completion of immediate execution.
        
        The behaviour of the {\bf dsp} attribute depends on the FPGA -
        \blist
        \item   For Virtex-5 and Virtex-6 it has two effects -
                \blist
                \item   it causes the static variable to be implemented using
                        the output register of a DSP block instead of
                        flip-flops.
                \item   For Virtex-5 and Virtex-6 it has a similar effect to
                        the {\bf alu} attribute in that it draws arithmetic and
                        bitwise logic operations in assignments to the
                        variable into a single ALU, but using a DSP block
                        rather than CLBs.
                \blend
        \item   For Virtex-2, Virtex-2Pro and Virtex-4\footnote{Virtex-4 has a
                DSP block but for reasons of local (non)usage the full
                capabilities of Virtex4 are not supported by 3PL. The DSP
                block is treated as a simple multiplier block.} it causes the
                static variable to be implemented using the output register
                of a multiplier block instead of flip-flops.
        \item   For other FPGAs the attribute has no effect.
        \blend
        
        The {\bf dsp} attribute will not implement the static variable using
        a multiplier or DSP register in the following cases -
        \blist
        \item   the variable is initialised to a non-zero value
        \item   the variable has the {\bf iob} attribute
        \blend
        The attribute will be ignored and a warning message issued in these
        cases.
        
        Note that the use of the {\bf dsp} attribute reduces the number of
        flip-flops required by using multiplier or DSP block output registers
        instead, but this may make timing constraints harder to meet because -
        \blist
        \item   the clock edge to output time of these may be worse that that
                for CLB flip-flops
        \item   the bits can no longer be distributed among CLBs and routing
                must instead converge on one or more multiplier or DSP blocks
        \blend
   
        In addition the {\bf loc} attribute if supplied will be ignored (with a
        warning message). For Virtex-2 and Virtex-2 PRO each multiplier element
        has a 36 bit register. For Virtex-5 and Virtex-6 the output register is
        48 bits. For variables whose type is not primitive each component will
        be implemented in a separate multiplier or DSP register. For components
        larger than the hardware register size sufficient DSP registers will be
        used to accomodate the component (but as noted above registers larger
        than 48 bits cannot use the Virtex-5 or Virtex-6 DSP block as an ALU).
        In the case of Virtex-5 or Virtex-6 , setting this attribute true also
        sets the {\bf alu} attribute true and the ALU logic will be implemented in the
        same DSP block. Note that if the limitations listed above are not met
        the {\bf dsp} attribute will be cleared (an associated warning message will be
        printed), the variable will be implemented using CLb flip-flops and if
        an ALU is generated it will use CLB logic.
        
        Note that for Virtex-5 or Virtex-6 if the {\bf alu} attribute is true but not
        the {\bf dsp} attribute then the static variable will be implemented using
        flip-flops and arithmetic operators '+' and '-' as the final or only
        operation can be combined into a single common adder/subtracter using
        CLBs.
        
        {\bf loc} provides absolute locations for the CLB flip-flops comprising
        the static variable (register). Locations are matched to flip-flops
        starting at the least significant bit. For a single bit static this
        attribute may be a string, otherwise it must be a string array. Each
        string must have the location format expected by the vendor for the FPGA
        type. 3PL does not at present provide support for relative placement
        of flip-flops in static variables though this could easily be added
        if required.

        {\bf iob} specifies that the flip-flops implementing the static variable
        are to be placed in IOB blocks rather than CLB blocks (this can also be
        requested by using the -pr option in the Xilinx map utility). If the
        value of the variable is used in addition to being connected to
        output buffers then the flip-flops will be duplicated in parallel, one
        set in the I/O blocks and the other in CLBs.
        
        {\bf ids} is a string array that contains the identifiers of the netlist
        elements (in this case D-type flip-flops) that implement the static mode
        variable (register). This attribute is read-only and is initialised when
        the variable is created. It can be used to apply constraints, such as
        absolute or relative location, to these elements.The array is ordered from
        least to most significant bit. For example to place a static variable in
        specific FPGA slice locations (for ISE) -
        
        \begin{lstlisting}[gobble=8, language=threepl]
        static(s, "uint:8");
        var(slocs, "[8]str", {"", "", "", .......}); // array of slice locations
        var(sids, "[8]str", attribute(s, ids));      // array of netlist element identifiers
        for (int(i) ; i<8 ; i++)
            addncfline("INST \""+ sids[i] + "\" LOC=" + slocs[i]);
        \end{lstlisting}
    
        For Vivado -
        
        \begin{lstlisting}[gobble=8, language=threepl]
        static(s, "uint:8");
        var(slocs, "[8]str", {"", "", "", .......}); // array of slice locations
        var(sids, "[8]str", attribute(s, ids));      // array of netlist element identifiers
        for (int(i) ; i<8 ; i++)
            addxdcline("set_property LOC " + slocs[i] + "[ get_cells " + sids[i] + "]");
        \end{lstlisting}
        
        
        
